\section{Pointwise norms: what is missing}\label{sec:pointwise-weak}

The covariance, the coarse Hessian and the determinant enter the step through $L^2$ bounds, and for them Section~\ref{sec:nested-walks} is complete. The background field is different. To localize the interaction and the composed activities, the weakened background map $A_*(s)$ must be built with weakened operators $\bH(s)$, $\bC(s)$, analytic on $\mathbb D_\kappa$. Its construction is a contraction in sup norms with derivatives (Proposition~\ref{prop:contraction}), or in $C^{1,\alpha}$ (blueprint, Appendix~C). So, for every $s\in\mathbb D_\kappa$, it needs bounds of the form
\begin{equation}\label{eq:pointwise-needed}
|\bC(s)f|_{1,\alpha;y}\le C\sum_{y'}e^{-\gamma|y-y'|}|f|_{B(y')},
\end{equation}
together with the same bounds for $\bH(s)$ and $\bC(s)D^T$.

\paragraph{Where the difficulty is.} The simplest way to obtain \eqref{eq:pointwise-needed} is to make the step operators $\mathcal R(s)$ bounded from sup norms to sup norms. Then every walk term is bounded term by term, as in \cite[(97)--(99)]{DimockI}. At level~(L3) this fails without further input, because of the seam term $h_j\one_{X'}(\mathfrak N(s)-\mathfrak N^{X'}(s))\one_{X'}G_j$.
\begin{itemize}
\item It is small only through the cancellation between its two parts. Weakening only one of them destroys the cancellation: for $|s_\Box|=e^\kappa$, the terms without steps are multiplied by $e^{\kappa\mathsf N_{\rm b}}$. So both parts are weakened, compatibly, and the terms of $\mathfrak N^{X'}(s)$ act on $u=G_jf$ throughout $\Lam^v_j$.
\item The interior estimate (Proposition~\ref{prop:system-interior}) bounds $u$ pointwise at distance at least $4$ from $\partial\Lam^v_j$; near $\partial\Lam^v_j$ only its $L^2$ norm is controlled.
\item The terms of $\mathfrak N^{X'}(s)$ that start near the boundary are products whose factors need pointwise information on their inputs. Interior regularity with $L^2$ sources controls values in $d=3$, but not gradients. The commutators $[\fA_{\rm loc},h_{j'}]$ are of first order.
\end{itemize}

Note, however, that absolute convergence does not force all factors to be sup-to-sup. It suffices that $\mathcal R(s)$ is a contraction in some Banach space adapted to the geometry. This leaves two routes.

\paragraph{Route 1: boundary regularity.} Prove regularity up to the boundary for the Dirichlet operators on the boxes of the construction:
\begin{quote}
\emph{(BR)} For a box $X$, the operators $G'_X$, $G'_XD^T_X$, $\bG^X$, $\bG^XD^T_X$, $\bH^X$ and $\bC^X$ satisfy the bounds of Theorem~\ref{thm:ID} and Proposition~\ref{prop:system-interior} at all blocks of $X$, not only at interior ones.
\end{quote}
Then all factors are sup-to-sup, and the construction of Section~\ref{sec:nested-walks} carries over to \eqref{eq:pointwise-needed}. Ba\l aban proves estimates of this kind for his Dirichlet propagators \cite{BalGreen,BalPropI,BalBackground}.

For the scalar Dirichlet Laplacian on a lattice box, odd reflection across the plane of the exterior layer gives a problem on the full lattice. At edges and corners, the reflections across the faces commute. The absorption argument of Section~\ref{sec:interior} can then be repeated on half-cubes and quarter-cubes at the boundary:
\begin{itemize}
\item the small perturbations become sources for the reflected equation. Their reflected coefficients may jump across the reflection plane, which is harmless only because $u$ vanishes on the exterior layer;
\item the averaging terms are block-local and the boxes are unions of blocks, so they remain $L^2$-controlled sources;
\item the mass term is of order zero.
\end{itemize}
The gauge case is not covered by this sketch. On $\ell^2_1(X)$ the bonds crossing $\partial X$ are kept, and since the exterior divergence is part of the operator, all components satisfy Dirichlet conditions for $D^TD+DD^T$ (Section~\ref{sec:reg-gauge}). The reflection then has to respect the lattice structure of the crossing bonds. $L^2$ regularity of second order on boxes is classical in the continuum, but the pointwise statement on the lattice, with these operators, has not been carried out.

\paragraph{Route 2: a mixed norm, without boundary estimates.} At the top level $X=\Tf$, every pointwise output is needed only on the supports $K_j$. Every Dirichlet cut involved lies at distance at least $\frac14(\rho-1)M$ from these supports, by \eqref{eq:rho-choice}. Two observations suggest that a norm mixing $L^2$ and pointwise control suffices:
\begin{itemize}
\item pieces of a kernel that jump a distance of at least $2$ need only $L^2$ input, by interior regularity on a cube where the source vanishes;
\item a chain of walk factors from a cut to $K_j$ passes through many interior steps. Each step of the form commutator times propagator gains one derivative, and each step of order zero times propagator gains two. So in $d=3$ a few interior steps take block-$L^2$ data to $C^{1,\alpha}$.
\end{itemize}
The norm would be block-$L^2$ within a few units of the cuts and pointwise elsewhere, with exponential weights. Its analytic input is a set of interior estimates with $L^p$ data, variants of Proposition~\ref{prop:system-interior}, together with the off-diagonal smoothing just described. These estimates are proved in Section~\ref{sec:interior-lp} (Corollary~\ref{cor:lp-dirichlet}). No boundary regularity is needed. The step bounds are carried out in Section~\ref{sec:tiers}, in block $L^6$ and H\"older norms rather than a mixed norm. They give \eqref{eq:pointwise-needed} for $\bG(s)$, $\bH(s)$, $\bC(s)$ and $\bC(s)D^T$ on the torus (Theorems~\ref{thm:tier-walk} and~\ref{thm:HC-pointwise}). A trace inequality on boxes (Lemma~\ref{lem:trace}) handles the Dirichlet cuts. $\bC(s)$ is represented through the torus identity $\bC=\bG\Pi^T-\bH\bQ\bG$.

\paragraph{Summary.} Sections~\ref{sec:weakenable} and~\ref{sec:nested-walks} give volume-type localization of everything that enters the step through $L^2$ bounds: the fluctuation covariance and its square root, the coarse Hessian, and the determinant. The background map, and with it the interaction and the composed activities, need in addition pointwise control of the weakened operators, \eqref{eq:pointwise-needed}. Either route above would provide it. Route~2 stays closer to the interior methods of this note. Route~1 is Ba\l aban's.
